\SourceLeaf{312}
Denique si \(a=0\), erit
\[
f=a'x'x'+2bx'x''+a''x''x''+2x(b''x'+b'x'')
\]
Accipiendo itaque \(x'\), \(x''\) ad libitum, ita tamen ut \(b''x'+b'x''\) non sit \(=0\) (quod manifesto fieri poterit, nisi simul \(b'\) et \(b''\) sint \(=0\); tunc autem foret \(D=0\)), nullo negotio perspicitur, \(x\) ita determinari posse, ut \(f\) obtineat valores tum positivos, tum negativos. Quare etiam in hocce casu \(f\) erit forma indefinita.

Eodem modo, ut hic ex numeris \(aD\), \(A'\) indolem formae \(f\) diiudicavimus, etiam \(aD\) et \(A''\) adhiberi possunt, ita ut \(f\) sit forma definita, si tum \(aD\) tum \(A''\) sit negativus; indefinita in omnibus reliquis casibus. Nec non prorsus simili modo eidem fini inservire potest consideratio numerorum \(a'D\) et \(A\), vel horum \(a'D\) et \(A''\), vel horum \(a''D\) et \(A\), vel denique ipsorum \(a''D\) et \(A'\).

Ex his omnibus colligitur, in forma definita sex numeros \(A\), \(A'\), \(A''\), \(aD\), \(a'D\), \(a''D\) esse negativos, et quidem in forma positiva \(a\), \(a'\), \(a''\) erunt positivi, \(D\) negativus; in negativa autem \(a\), \(a'\), \(a''\) erunt negativi, \(D\) positivus. Hinc patet, omnes formas ternarias determinantis dati positivi distribui in negativas et indefinitas; omnes autem determinantis negativi in positivas et indefinitas; denique formas positivas determinantis positivi, seu negativas determinantis negativi omnino non dari. --- Indidem facile perspicitur, formae definitae semper adiunctam esse definitam et quidem \emph{negativam}, indefinitae indefinitam.

Quum omnes numeri per formam ternariam datam repraesentabiles manifesto etiam per omnes formas huic aequivalentes repraesentari possint: formae ternariae in eadem classe contentae vel omnes erunt indefinitae, vel omnes positivae, vel omnes negativae. Quamobrem has formarum denominationes etiam ad classes integras transferre licebit.

\Needspace{8\baselineskip}
\Article{272}
Theorema in art. praec. propositum, quod omnes formae ternariae determinantis dati in multitudinem \emph{finitam} classium distribuuntur, per methodum ei qua in formis binariis usi sumus analogam tractabimus, scilicet ostendendo, primo, quo pacto quaevis forma ternaria ad formam simpliciorem reduci possit, dein, formarum simplicissimarum (ad quas per tales reductiones perveniatur) multitudinem pro quovis determinante dato esse finitam. Supponamus generali%
\SourceLeaf{313}%
ter, propositam esse formam ternariam \(f=\left(\begin{smallmatrix}a,&a',&a''\\b,&b',&b''\end{smallmatrix}\right)\) determinantis \(D\) (a cifra diversi), quae per substitutionem \((S)\)
\[
\begin{array}{ccc}
\alpha,&\beta,&\gamma\\
\alpha',&\beta',&\gamma'\\
\alpha'',&\beta'',&\gamma''
\end{array}
\]
transeat in aequivalentem \(g=\left(\begin{smallmatrix}m,&m',&m''\\n,&n',&n''\end{smallmatrix}\right)\); versabiturque negotium nostrum in eo, ut \(\alpha\), \(\beta\), \(\gamma\) etc. ita definiantur, ut forma \(g\) simplicior evadat quam \(f\). Sint formae ipsis \(f\), \(g\) adiunctae resp. \(\left(\begin{smallmatrix}A,&A',&A''\\B,&B',&B''\end{smallmatrix}\right)\), \(\left(\begin{smallmatrix}M,&M',&M''\\N,&N',&N''\end{smallmatrix}\right)\). quae designentur per \(F\), \(G\). Tunc per art. 269 \(F\) transibit in \(G\) per substitutionem ipsi \(S\) adiunctam, \(G\) autem in \(F\) per substitutionem ex transpositione ipsius \(S\) oriundam. Numerum
\[
\alpha\beta'\gamma''+\alpha'\beta''\gamma+\alpha''\beta\gamma'-\alpha''\beta'\gamma-\alpha\beta''\gamma'-\alpha'\beta\gamma''
\]
qui esse debebit vel \(=+1\) vel \(=-1\), denotabimus per \(k\). Quibus ita factis observamus

I. Si fiat \(\gamma=0\), \(\gamma'=0\), \(\alpha''=0\), \(\beta''=0\), \(\gamma''=1\), fore
\[
\begin{aligned}
m&=a\alpha\alpha+2b''\alpha\alpha'+a'\alpha'\alpha',\\
m'&=a\beta\beta+2b''\beta\beta'+a'\beta'\beta',\\
m''&=a''\\
n&=b\beta'+b'\beta,\\
n'&=b\alpha'+b'\alpha,\\
n''&=a\alpha\beta+b''(\alpha\beta'+\beta\alpha')+a'\alpha'\beta'
\end{aligned}
\]
Praeterea esse debebit \(\alpha\beta'-\beta\alpha'\) vel \(=+1\) vel \(=-1\). Hinc manifestum est, formam binariam \((a,b'',a')\), cuius determinans est \(A''\), transmutari per substitutionem \(\alpha\), \(\beta\), \(\alpha'\), \(\beta'\) in formam binariam \((m,n'',m')\) determinantis \(M''\), et proin ipsi aequivalere propter \(\alpha\beta'-\beta\alpha'=\pm1\), unde erit \(M''=A''\), quod etiam directe facile confirmatur. Nisi itaque \((a,b'',a')\) iam est forma simplicissima in classe sua, ipsos \(\alpha\), \(\beta\), \(\alpha'\), \(\beta'\) ita determinare licebit, ut \((m,n'',m')\) sit forma simplicior; et quidem e theoria aequivalentiae formarum binariarum facile concluditur, hoc ita fieri posse, ut \(m\) non sit maior quam \(\sqrt{-\frac{4}{3}A''}\), si \(A''\) fuerit negativus, vel non maior quam \(\sqrt{A''}\), si \(A''\) fuerit positivus, vel \(m=0\), si \(A''=0\), ita ut in omnibus casibus valor (absolutus) ipsius \(m\) certe vel infra vel saltem usque ad \(\sqrt{\pm\frac{4}{3}A''}\) deprimi possit. Hoc itaque modo forma \(f\) ad aliam reducitur co\"efficientem primum, si fieri potest, minorem habentem, et cuius forma adiuncta co\"efficientem tertium eundem habet ut forma \(F\) ipsi \(f\) adiuncta. In hoc consistit \emph{reductio prima}.

\SourceLeaf{314}
II. Si vero fit \(\alpha=1\), \(\beta=0\), \(\gamma=0\), \(\alpha'=0\), \(\alpha''=0\), erit \(k=\beta'\gamma''-\beta''\gamma'=\pm1\); substitutio itaque ipsi \(S\) adiuncta erit
\[
\begin{array}{rrr}
\pm1,&0,&0\\
0,&\gamma'',&-\beta''\\
0,&-\gamma',&\beta'
\end{array}
\]
per quam \(F\) transibit in \(G\). Habebitur itaque
\[
\begin{aligned}
m&=a,\quad n'=b'\gamma''+b''\gamma',\quad n''=b'\beta''+b''\beta'\\
m'&=a'\beta'\beta'+2b\beta'\beta''+a''\beta''\beta''\\
m''&=a'\gamma'\gamma'+2b\gamma'\gamma''+a''\gamma''\gamma''\\
n&=a'\beta'\gamma'+b(\beta'\gamma''+\gamma'\beta'')+a''\beta''\gamma''\\
M'&=A'\gamma''\gamma''-2B\gamma'\gamma''+A''\gamma'\gamma'\\
N&=-A'\beta''\gamma''+B(\beta'\gamma''+\gamma'\beta'')-A''\beta'\gamma'\\
M''&=A'\beta''\beta''-2B\beta'\beta''+A''\beta'\beta'
\end{aligned}
\]
Hinc patet, formam binariam \((A'',B,A')\), cuius determinans est \(Da\), transire per substitutionem \(\beta'\), \(-\gamma'\), \(-\beta''\), \(\gamma''\) in formam \((M'',N,M')\) determinantis \(Dm\), adeoque (propter \(\beta'\gamma''-\gamma'\beta''=\pm1\), vel propter \(Da=Dm\)) ipsi aequivalere. Nisi itaque \((A'',B,A')\) iam est forma simplicissima classis suae, co\"efficientes \(\beta'\), \(\gamma'\), \(\beta''\), \(\gamma''\) ita determinari poterunt, ut \((M''.\,N,M')\) sit simplicior, et quidem hoc semper poterit fieri ita, ut \(M''\) sine respectu signi non sit maior quam \(\sqrt{\pm\frac{4}{3}Da}\). Hoc itaque modo forma \(f\) reducitur ad aliam co\"efficientem primum eundem habentem, sed cuius forma adiuncta co\"efficientem tertium si fieri potest minorem habeat quam forma \(F\) ipsi \(f\) adiuncta. In hoc consistit \emph{reductio secunda}.

III. Si itaque \(f\) est forma ternaria, ad quam neque reductio prima neque secunda est applicabilis, \emph{i. e.} quae per neutram in formam simpliciorem transmutari potest: necessario erit tum \(aa<\) vel \(=\frac{4}{3}A''\), tum \(A''A''<\) vel \(=\frac{4}{3}aD\) sine respectu signi. Hinc \(a^4\) erit \(<\) vel \(=\frac{16}{9}A''A''\), adeoque \(a^4<\) vel \(=\frac{64}{27}aD\), \(a^3<\) vel \(=\frac{64}{27}D\), et \(a<\) vel \(=\frac{4}{3}\sqrt[3]{D}\); hinc rursus \(A''A''<\) vel \(=\frac{16}{9}\sqrt[3]{D^4}\) atque \(A''<\) vel \(=\frac{4}{3}\sqrt[3]{D^2}\). Quamobrem quamdiu \(a\) vel \(A''\) hos limites adhuc superant, necessario una aut altera reductionum praecedentium ad formam \(f\) applicari poterit. --- Ceterum haec conclusio non est convertenda, quum utique saepius
\SourceLeaf{315}
accidat, ut forma ternaria, cuius co\"efficiens primus, atque co\"efficiens tertius formae adiunctae iam sunt infra illos limites, nihilominus per unam alteramve reductionem adhuc simplicior reddi possit.

IV. Quodsi vero ad formam ternariam quamcunque datam determinantis \(D\) alternis vicibus reductio prima et secunda applicantur, \emph{i. e.} ad ipsam prima vel secunda, ad eam quae hinc resultat secunda vel prima, ad eam quae hinc provenit iterum prima vel secunda etc., manifestum est, tandem necessario ad formam perventum iri, ad quam neutra amplius applicari possit. Quum enim magnitudo absoluta tum co\"efficientium primorum formarum hoc modo prodeuntium, tum co\"efficientium tertiorum formarum illis adiunctarum continuo alternis vicibus eadem maneat atque decrescat, hic progressus necessario tandem alicubi finietur, quia alioquin duae series infinitae numerorum continuo decrescentium haberentur. Hinc iam nacti sumus egregium theorema: \emph{Quaevis forma ternaria determinantis \(D\) reduci potest ad aliam aequivalentem, cuius co\"efficiens primus non sit maior quam \(\frac{4}{3}\sqrt[3]{D}\), atque co\"efficiens tertius formae ipsi adiunctae non maior quam \(\frac{4}{3}\sqrt[3]{D^2}\) sine respectu signi, siquidem forma proposita his proprietatibus ipsa nondum est praedita.} --- Ceterum loco co\"efficientis primi formae \(f\) atque tertii formae ipsi \(f\) adiunctae prorsus simili modo tractare potuissemus vel co\"efficientem primum formae ipsius et secundum adiunctae; vel secundum formae ipsius et primum vel tertium adiunctae; vel tertium formae ipsius et primum vel secundum adiunctae, quibus viis perinde ad finem nobis propositum perveniremus: sed e re est, methodo uni constanter adhaerere, quo facilius operationes huc pertinentes ad algorithmum fixum reduci possint. Denique observamus, duobus co\"efficientibus, quos infra limites fixos deprimere docuimus. limites adhuc minores constitui posse, si formae definitae ab indefinitis separentur; hoc vero ad institutum praesens non est necessarium.

\Needspace{8\baselineskip}
\Article{273}
Ecce iam quaedam exempla, per quae praecepta praecedentia magis illustrabuntur.

\emph{Ex.} 1. Sit \(f=\left(\begin{smallmatrix}19,&21,&50\\15,&28,&1\end{smallmatrix}\right)\), erit \(F=\left(\begin{smallmatrix}-825,&-166,&-398\\257,&573,&-370\end{smallmatrix}\right)\), \(D=-1\). Quum \((19,1,21)\) sit forma binaria reducta, cui alia, termini primi minoris quam \(19\), non aequivalet, reductio prima hic non est applicabilis; forma binaria
\SourceLeaf{316}
\((A'',B,A')=(-398,257,-166)\) autem per theoriam aequivalentiae formarum binariarum in simpliciorem aequivalentem \((-2,1,-10)\) transmutabilis invenitur, in quam transit per substitutionem \(2,7,3,11\). Faciendo itaque \(\beta'=2\), \(\gamma'=-7\), \(\beta''=-3\), \(\gamma''=11\), applicanda erit ad formam \(f\) substitutio \(\left\{\begin{smallmatrix}1,&0,&0\\0,&2,&-7\\0,&-3,&11\end{smallmatrix}\right\}\) per quam invenitur transire in hanc \(\left(\begin{smallmatrix}19,&354,&4769\\-1299,&301,&-82\end{smallmatrix}\right)\ldots f'\). Co\"efficiens tertius formae, huic adiunctae, est \(-2\), quo respectu \(f'\) simplicior est censenda quam \(f\).

Ad formam \(f'\) applicari potest reductio prima. Scilicet quum forma binaria \((19,-82,354)\) transmutetur in \((1,0,2)\) per substitutionem \(13,4,3,1\): applicanda erit ad formam \(f'\) substitutio \(\left\{\begin{smallmatrix}13,&4,&0\\3,&1,&0\\0,&0,&1\end{smallmatrix}\right\}\) per quam transit in hanc \(\left(\begin{smallmatrix}1,&2,&4769\\-95,&16,&0\end{smallmatrix}\right)\ldots f''\).

Ad formam \(f''\), cui adiuncta est \(\left(\begin{smallmatrix}-513,&-4513,&-2\\-95,&32,&1520\end{smallmatrix}\right)\), denuo applicari potest reductio secunda. Scilicet \((-2,-95,-4513)\) transit per substitutionem \(47,1,-1,0\) in \((-1,1,-2)\): quamobrem ad \(f''\) applicanda erit substitutio \(\left\{\begin{smallmatrix}1,&0,&0\\0,&47,&-1\\0,&1,&0\end{smallmatrix}\right\}\) per quam transit in \(\left(\begin{smallmatrix}1,&257,&2\\1,&0,&16\end{smallmatrix}\right)\ldots f'''\). Huius co\"efficiens primus per reductionem primam amplius diminui non potest, neque formae, ipsi adiunctae, tertius per secundam.

\emph{Ex.} 2. Proposita sit forma \(\left(\begin{smallmatrix}10,&26,&2\\7,&0,&4\end{smallmatrix}\right)\ldots f\), cui adiuncta est \(\left(\begin{smallmatrix}-3,&-20,&-244\\70,&-28,&8\end{smallmatrix}\right)\) et cuius determinans \(=2\). Hic successive reperiuntur, applicando alternatim reductionem secundam et primam,
\[
\begin{array}{c|c|c}
\text{substitutiones}&\text{per quas transit}&\text{in}\\[3pt]
\left\{\begin{smallmatrix}1,&0,&0\\0,&-1,&0\\0,&4,&-1\end{smallmatrix}\right\}
&f&\left(\begin{smallmatrix}10,&2,&2\\-1,&0,&-4\end{smallmatrix}\right)=f'\\[6pt]
\left\{\begin{smallmatrix}0,&-1,&0\\1,&-2,&0\\0,&0,&1\end{smallmatrix}\right\}
&f'&\left(\begin{smallmatrix}2,&2,&2\\2,&-1,&0\end{smallmatrix}\right)=f''\\[6pt]
\left\{\begin{smallmatrix}1,&0,&0\\0,&-1,&0\\0,&2,&-1\end{smallmatrix}\right\}
&f''&\left(\begin{smallmatrix}2,&2,&2\\-2,&1,&-2\end{smallmatrix}\right)=f'''\\[6pt]
\left\{\begin{smallmatrix}1,&0,&0\\1,&1,&0\\0,&0,&1\end{smallmatrix}\right\}
&f'''&\left(\begin{smallmatrix}0,&2,&2\\-2,&-1,&0\end{smallmatrix}\right)=f''''
\end{array}
\]
Forma \(f''''\) per reductionem primam vel secundam ulterius deprimi nequit.

\Needspace{8\baselineskip}
\Article{274}
Quando forma ternaria habetur, cuius co\"efficiens primus, atque formae adiunctae tertius, quantum fieri potest per methodos praecedentes sunt depressi: methodus sequens reductionem ulteriorem suppeditat.

\SourceLeaf{317}
Adhibendo signa eadem ut in art. 272, et ponendo \(\alpha=1\), \(\alpha'=0\), \(\beta'=1\), \(\alpha''=0\), \(\beta''=0\), \(\gamma''=1\), \emph{i. e.} adhibendo substitutionem
\[
\begin{array}{ccc}
1,&\beta,&\gamma\\
0,&1,&\gamma'\\
0,&0,&1
\end{array}
\]
erit
\[
\begin{aligned}
m&=a,\quad m'=a'+2b''\beta+a\beta\beta,\\
m''&=a''+2b\gamma'+2b'\gamma+a\gamma\gamma+2b''\gamma\gamma'+a'\gamma'\gamma'\\
n&=b+a'\gamma'+b'\beta+b''(\gamma+\beta\gamma')+a\beta\gamma,\\
n'&=b'+a\gamma+b''\gamma',\quad n''=b''+a\beta
\end{aligned}
\]
praeterea
\[
M''=A'',\quad N=B-A''\gamma',\quad N'=B'-N\beta-A''\gamma
\]
Per talem itaque substitutionem co\"efficientes \(a\), \(A''\), qui per reductiones praecedentes diminuti sunt, non mutantur; quamobrem negotium in eo versatur, ut per idoneam determinationem ipsorum \(\beta\), \(\gamma\), \(\gamma'\) depressiones in co\"efficientibus reliquis obtineantur. Ad hunc finem observamus primo, si fuerit \(A''=0\), supponi posse, esse etiam \(a=0\); si enim \(a\) non \(=0\), reductio prima adhuc semel applicabilis foret, quum cuivis formae binariae determinantis \(0\) aequivaleat forma talis \((0,0,h)\), sive cuius terminus primus \(=0\) (V. art. 215). Prorsus simili ratione supponere licet, esse etiam \(A''=0\), si fuerit \(a=0\), ita ut vel neuter numerorum \(a\), \(A''\) sit \(0\) vel uterque.

In casu priori manifestum est, ipsos \(\beta\), \(\gamma\), \(\gamma'\) ita determinari posse, ut sine respectu signi \(n''\), \(N\), \(N'\) resp. non sint maiores quam \(\frac{1}{2}a\), \(\frac{1}{2}A''\), \(\frac{1}{2}A''\). Ita in exemplo primo art. praec. transibit forma postrema \(\left(\begin{smallmatrix}1,&257,&2\\1,&0,&16\end{smallmatrix}\right)\), cui adiuncta est \(\left(\begin{smallmatrix}-513,&-2,&-1\\1,&-16,&32\end{smallmatrix}\right)\), per substitutionem \(\left\{\begin{smallmatrix}1,&-16,&16\\0,&1,&-1\\0,&0,&1\end{smallmatrix}\right\}\) in hanc \(\left(\begin{smallmatrix}1,&1,&1\\0,&0,&0\end{smallmatrix}\right)\ldots f''''\), cui adiuncta est \(\left(\begin{smallmatrix}-1,&-1,&-1\\0,&0,&0\end{smallmatrix}\right)\).

In casu posteriori, ubi \(a=A''=0\), adeoque etiam \(b''=0\) erit
\[
\begin{aligned}
m&=0,\quad m'=a',\quad m''=a''+2b\gamma'+2b'\gamma+a'\gamma'\gamma'\\
n&=b+a'\gamma'+b'\beta,\quad n'=b',\quad n''=0
\end{aligned}
\]
Erit itaque
\[
D=a'b'b'=m'n'n'
\]
perspicieturque facile, \(\beta\) et \(\gamma'\) ita determinari posse, ut \(n\) fiat aequalis residuo absolute minimo ipsius \(b\) secundum modulum, qui est divisor communis maximus
\SourceLeaf{318}
ipsorum \(a'\), \(b'\); i. e. ut \(n\) fiat non maior quam semissis huius divisoris sine respectu signi, adeoque \(n=0\), quoties \(a'\), \(b'\) inter se sunt primi. Ipsis \(\beta\), \(\gamma'\) in hunc modum determinatis, valor ipsius \(\gamma\) ita accipi poterit, ut \(m''\) non sit maior quam \(b'\) sine respectu signi; hoc quidem impossibile esset, quando \(b'=0\); tunc vero foret \(D=0\), quem casum exclusimus. Ita fit pro forma postrema in \emph{ex.} 2 art. praec. \(n=-2-\beta+2\gamma'\), unde statuendo \(\beta=-2\), \(\gamma'=0\), fit \(n=0\), porro \(m''=2-2\gamma\), et ponendo \(\gamma=1\), \(m''=0\). Habemus itaque substitutionem \(\left\{\begin{smallmatrix}1,&-2,&1\\0,&1,&0\\0,&0,&1\end{smallmatrix}\right\}\) per quam forma illa transit in \(\left(\begin{smallmatrix}0,&2,&0\\0,&-1,&0\end{smallmatrix}\right)\ldots f'''''\).

\Needspace{8\baselineskip}
\Article{275}
Si habetur series formarum ternariarum aequivalentium \(f\), \(f'\), \(f''\), \(f'''\) etc., atque transformationes cuiusvis harum formarum in sequentem: ex transformationibus formae \(f\) in \(f'\), formaeque \(f'\) in \(f''\) per art. 270 deducitur transformatio formae \(f\) in \(f''\); ex hac atque transf. formae \(f''\) in \(f'''\) sequitur transf. formae \(f\) in \(f'''\) etc., manifestoque hoc pacto transformatio formae \(f\) in quamcunque aliam seriei inveniri poterit. Et quum ex transformatione formae \(f\) in quamcunque aliam aequivalentem \(g\) deduci possit transformatio formae \(g\) in \(f\) (\(S''\) ex \(S\) artt. 268, 269), hoc modo erui poterit transformatio cuiuslibet formae seriei \(f'\), \(f''\) etc. in primam \(f\). --- Ita pro formis exempli primi art. praec. inveniuntur substitutiones
\[
\begin{array}{rrr|rrr|rrr}
13,&4,&0&13,&188,&-4&13,&-20,&16\\
6,&2,&-7&6,&87,&-2&6,&-9,&7\\
-9,&-3,&11&-9,&-130,&3&-9,&14,&-11
\end{array}
\]
per quas \(f\) transit in \(f''\), \(f'''\), \(f''''\) resp., et ex subst. ultima haec \(\left\{\begin{smallmatrix}1,&4,&4\\3,&1,&5\\3,&-2,&3\end{smallmatrix}\right\}\) per quam \(f''''\) transit in \(f\). Simili modo pro \emph{ex.} 2 art. praec. prodeunt substitutiones
\[
\begin{array}{rrr|rrr}
1,&-1,&1&2,&-3,&-1\\
-3,&4,&-3&3,&1,&0\\
10,&-14,&11&2,&4,&1
\end{array}
\]
per quas resp. transit forma \(\left(\begin{smallmatrix}10,&26,&2\\7,&0,&4\end{smallmatrix}\right)\) in \(\left(\begin{smallmatrix}0,&2,&0\\0,&-1,&0\end{smallmatrix}\right)\), atque haec in illam.

\Needspace{8\baselineskip}
\Article{276}
\textsc{Theorema.} \emph{Classium, in quas omnes formae ternariae determinantis dati distribuuntur, multitudo semper est finita.}

\SourceLeaf{319}
\emph{Dem.} I. Multitudo omnium formarum \(\left(\begin{smallmatrix}a,&a',&a''\\b,&b',&b''\end{smallmatrix}\right)\) determinantis dati \(D\), in quibus \(a=0\), \(b''=0\), \(b\) non maior quam semissis divisoris comm. max. numerorum \(a'\), \(b'\); \(a''\) non maior quam \(b'\), manifesto est finita. Quoniam enim esse debet \(a'b'b'=D\), pro \(b'\) alii valores accipi nequeunt, quam \(+1\), \(-1\) atque radices quadratorum ipsum \(D\) metientium (si quae alia praeter \(1\) dantur) signo positivo et negativo affectae, quorum valorum multitudo finita est. Pro singulis autem valoribus ipsius \(b'\) valor ipsius \(a'\) est determinatus, ipsorumque \(b\), \(a''\) valores manifesto limitantur ad multitudinem finitam.

II. Simili modo finita est multitudo omnium formarum \(\left(\begin{smallmatrix}a,&a',&a''\\b,&b',&b''\end{smallmatrix}\right)\) determinantis \(D\), in quibus \(a\) non \(=0\), neque maior quam \(\frac{4}{3}\sqrt[3]{\pm D}\); \(b''b''-aa'=A''\) non \(=0\) neque maior quam \(\frac{4}{3}\sqrt[3]{D^2}\); \(b''\) non maior quam \(\frac{1}{2}a\); \(ab-b'b''=B\) et \(a'b'-bb''=B'\) non maiores quam \(\frac{1}{2}A''\). Nam multitudo omnium combinationum valorem ipsorum \(a\), \(b''\), \(A''\), \(B\), \(B'\) finita erit; his vero singulis determinatis, etiam formae co\"efficientes reliqui \(a'\), \(b\), \(b'\), \(a''\), co\"efficientesque formae adiunctae
\[
bb-a'a''=A,\quad b'b'-aa''=A',\quad a''b''-bb'=B''
\]
determinati erunt per aequationes hasce:
\[
\begin{aligned}
a'&=\frac{b''b''-A''}{a},\quad A'=\frac{BB-aD}{A''},\\
A&=\frac{B'B'-a'D}{A''},\quad B''=\frac{BB'+b''D}{A''}\\
b&=\frac{AB-B'B''}{D}=-\frac{Ba'+B'b''}{A''},\\
b'&=\frac{A'B'-BB''}{D}=-\frac{Bb''+B'a}{A''}\\
a''&=\frac{b'b'-A'}{a}=\frac{bb-A}{a'}=\frac{bb'+B''}{b''}
\end{aligned}
\]
Iam quum omnes illae formae obtineantur, eligendo e cunctis combinationibus valorum ipsorum \(a\), \(b''\), \(A''\), \(B\), \(B'\) eas, e quibus etiam \(a'\), \(a''\), \(b\), \(b'\) valores integros nanciscuntur, illarum multitudo manifesto erit finita.

III. Cunctae itaque formae in I et II multitudinem finitam classium constituunt, quae etiam formarum ipsarum multitudine minor esse poterit, si quae ex ipsis inter se sunt aequivalentes. Iam quum per disquisitiones praecedentes quaevis forma ternaria determinantis \(D\) alicui ex illis formis necessario aequivaleat, \emph{i. e.} ad aliquam e classibus, quas hae formae constituunt, pertineat: hae classes omnes formas det. \(D\) complectentur, \emph{i. e.} omnes formae ternariae det. \(D\) in multitudinem finitam classium distribuentur. \emph{Q. E. D.}

\Needspace{8\baselineskip}
\SourceLeaf{320}
\Needspace{8\baselineskip}
\Article{277}
Regulae, per quas omnes formae in I et II art. praec. erui possunt, ex ipsarum explicatione sponte defluunt; quare sufficiet quaedam exempla apposuisse. Pro \(D=1\), formae I hae sex (per ambiguitatem signorum) prodeunt.
\[
\left(\begin{smallmatrix}0,&1,&0\\0,&\pm1,&0\end{smallmatrix}\right),\quad
\left(\begin{smallmatrix}0,&1,&\pm1\\0,&\pm1,&0\end{smallmatrix}\right)
\]
in formis II \(a\) et \(A''\) alios valores quam \(+1\) et \(-1\) habere nequeunt, pro singulis quatuor combinationum hinc oriundarum \(b''\), \(B\) et \(B'\) poni debent \(=0\), unde emergunt quatuor formae
\[
\left(\begin{smallmatrix}1,&-1,&1\\0,&0,&0\end{smallmatrix}\right),\quad
\left(\begin{smallmatrix}-1,&1,&1\\0,&0,&0\end{smallmatrix}\right),\quad
\left(\begin{smallmatrix}1,&1,&-1\\0,&0,&0\end{smallmatrix}\right),\quad
\left(\begin{smallmatrix}-1,&-1,&-1\\0,&0,&0\end{smallmatrix}\right)
\]
Simili modo pro \(D=-1\) sex formae I quatuorque II habentur,
\[
\begin{gathered}
\left(\begin{smallmatrix}0,&-1,&0\\0,&\pm1,&0\end{smallmatrix}\right),\quad
\left(\begin{smallmatrix}0,&-1,&\pm1\\0,&\pm1,&0\end{smallmatrix}\right);\quad
\left(\begin{smallmatrix}1,&-1,&-1\\0,&0,&0\end{smallmatrix}\right),\\
\left(\begin{smallmatrix}-1,&1,&-1\\0,&0,&0\end{smallmatrix}\right),\quad
\left(\begin{smallmatrix}-1,&-1,&1\\0,&0,&0\end{smallmatrix}\right),\quad
\left(\begin{smallmatrix}1,&1,&1\\0,&0,&0\end{smallmatrix}\right)
\end{gathered}
\]
Pro \(D=2\) sex formae I proveniunt
\[
\left(\begin{smallmatrix}0,&2,&0\\0,&\pm1,&0\end{smallmatrix}\right),\quad
\left(\begin{smallmatrix}0,&2,&\pm1\\0,&\pm1,&0\end{smallmatrix}\right)
\]
octoque formae II
\[
\begin{gathered}
\left(\begin{smallmatrix}1,&-1,&2\\0,&0,&0\end{smallmatrix}\right),\quad
\left(\begin{smallmatrix}-1,&1,&2\\0,&0,&0\end{smallmatrix}\right),\quad
\left(\begin{smallmatrix}1,&1,&-2\\0,&0,&0\end{smallmatrix}\right),\quad
\left(\begin{smallmatrix}-1,&-1,&-2\\0,&0,&0\end{smallmatrix}\right)\\
\left(\begin{smallmatrix}1,&-2,&1\\0,&0,&0\end{smallmatrix}\right),\quad
\left(\begin{smallmatrix}-1,&2,&1\\0,&0,&0\end{smallmatrix}\right),\quad
\left(\begin{smallmatrix}1,&2,&-1\\0,&0,&0\end{smallmatrix}\right),\quad
\left(\begin{smallmatrix}-1,&-2,&-1\\0,&0,&0\end{smallmatrix}\right)
\end{gathered}
\]
Ceterum multitudo classium ex his formis in his tribus casibus prodeuntium formarum multitudine multo minor est. Scilicet facile confirmatur

I. Formam \(\left(\begin{smallmatrix}0,&1,&0\\0,&1,&0\end{smallmatrix}\right)\) transire in
\[
\left(\begin{smallmatrix}0,&1,&0\\0,&-1,&0\end{smallmatrix}\right),\quad
\left(\begin{smallmatrix}0,&1,&1\\0,&\pm1,&0\end{smallmatrix}\right),\quad
\left(\begin{smallmatrix}0,&1,&-1\\0,&\pm1,&0\end{smallmatrix}\right),\quad
\left(\begin{smallmatrix}1,&1,&-1\\0,&0,&0\end{smallmatrix}\right)
\]
resp. per substitutiones
\[
\begin{array}{rrr|rrr|rrr|rrr}
1,&0,&0&0,&0,&1&0,&0,&1&1,&0,&-1\\
0,&1,&0&0,&1,&-1&0,&1,&1&1,&1,&-1\\
0,&0,&-1&\pm1,&1,&0&\pm1,&-1,&-1&0,&-1,&1
\end{array}
\]
formam \(\left(\begin{smallmatrix}1,&1,&-1\\0,&0,&0\end{smallmatrix}\right)\) autem in \(\left(\begin{smallmatrix}1,&-1,&1\\0,&0,&0\end{smallmatrix}\right)\), \(\left(\begin{smallmatrix}-1,&1,&1\\0,&0,&0\end{smallmatrix}\right)\) per solam indeterminatarum permutationem. Quare illae decem formae ternariae det. \(1\) ad has duas reducuntur
\SourceLeaf{321}
\(\left(\begin{smallmatrix}0,&1,&0\\0,&1,&0\end{smallmatrix}\right)\), \(\left(\begin{smallmatrix}-1,&-1,&-1\\0,&0,&0\end{smallmatrix}\right)\); pro priori, si magis arridet, etiam haec \(\left(\begin{smallmatrix}1,&0,&0\\1,&0,&0\end{smallmatrix}\right)\) accipi potest. Quum forma prior indefinita sit, posterior definita, manifestum est, quamvis formam ternariam indefinitam det. \(1\) aequivalere formae \(xx+2yz\), quamvis definitam huic \(-xx-yy-zz\).

II. Prorsus simili modo invenitur, quamlibet formam ternariam indefinitam determinantis \(-1\) aequivalere formae \(-xx+2yz\), quamlibet definitam huic \(xx+yy+zz\).

III. Pro determinante \(2\) ex octo formis (II) statim reiici possunt secunda, sexta et septima, quippe quae ex prima per solam indeterminatarum permutationem oriuntur, similique ratione etiam quinta quae e tertia, et octava quae e quarta perinde proveniunt; tres reliquae cum sex formis I, tres classes constituunt; scilicet \(\left(\begin{smallmatrix}0,&2,&0\\0,&1,&0\end{smallmatrix}\right)\) transit in \(\left(\begin{smallmatrix}0,&2,&0\\0,&-1,&0\end{smallmatrix}\right)\) per substitutionem \(\left\{\begin{smallmatrix}1,&0,&0\\0,&1,&0\\0,&0,&-1\end{smallmatrix}\right\}\) formaque \(\left(\begin{smallmatrix}1,&1,&-2\\0,&0,&0\end{smallmatrix}\right)\) in
\[
\begin{gathered}
\left(\begin{smallmatrix}0,&2,&1\\0,&1,&0\end{smallmatrix}\right),\quad
\left(\begin{smallmatrix}0,&2,&1\\0,&-1,&0\end{smallmatrix}\right),\quad
\left(\begin{smallmatrix}0,&2,&-1\\0,&1,&0\end{smallmatrix}\right),\\
\left(\begin{smallmatrix}0,&2,&-1\\0,&-1,&0\end{smallmatrix}\right),\quad
\left(\begin{smallmatrix}1,&-1,&2\\0,&0,&0\end{smallmatrix}\right)
\end{gathered}
\]
resp. per substitutiones
\[
\begin{array}{rrr|rrr|rrr|rrr|rrr}
1,&0,&1&1,&0,&-1&1,&0,&0&1,&0,&0&1,&0,&0\\
1,&2,&0&1,&2,&0&1,&2,&-1&1,&2,&1&0,&1,&2\\
1,&1,&0&1,&1,&0&1,&1,&-1&1,&1,&1&0,&1,&1
\end{array}
\]
Quaevis itaque forma ternaria determinantis \(2\) ad aliquam ex his tribus est reducibilis
\[
\left(\begin{smallmatrix}0,&2,&0\\0,&-1,&0\end{smallmatrix}\right),\quad
\left(\begin{smallmatrix}1,&1,&-2\\0,&0,&0\end{smallmatrix}\right),\quad
\left(\begin{smallmatrix}-1,&-1,&-2\\0,&0,&0\end{smallmatrix}\right)
\]
loco primae, si magis placet, etiam \(\left(\begin{smallmatrix}2,&0,&0\\1,&0,&0\end{smallmatrix}\right)\) accipi potest. Manifesto autem quaevis forma ternaria definita necessario aequivalebit tertiae \(-xx-yy-2zz\), quum duae priores sint indefinitae: quaevis indefinita primae vel secundae, et quidem primae \(2xx+2yz\), si ipsius co\"efficiens primus, secundus et tertius simul sunt pares (quoniam facile perspicitur, talem formam per substitutionem quamcunque in similem formam transire, adeoque formae secundae aequivalere non posse), secundae \(xx+yy-2zz\) autem, si ipsius co\"efficiens primus, secundus et tertius non simul pares sunt, sed unus, duo omnesve impares (in talem enim formam ex simili ratione forma prima \(2xx+2yz\) per nullam substitutionem transformabilis esse poterit).
